\chapter{Concrete Examples}\label{ch:examples}

This chapter collects concrete MPS tensors that illustrate the general
theory developed in earlier chapters.  Each example is specified by its
physical dimension~$d$, bond dimension~$D$, and explicit matrix
family~$\{A^i\}$.  Key structural properties (injectivity, RFP status,
symmetry) are stated as propositions linking back to the general
definitions.

References follow~\cite{Cirac2021Matrix}, Section~III.

%% ===================================================================
\section{The GHZ state}\label{sec:ghz}
%% ===================================================================

The Greenberger--Horne--Zeilinger (GHZ) state is the prototypical
non-injective, renormalization-fixed-point MPS.

\begin{definition}[GHZ tensor]\label{def:ghz_tensor}
    \lean{MPSTensor.ghzTensor}
    \leanok
    \uses{def:mps_tensor}
    Set $d = D = 2$.  The GHZ tensor is the family
    $\{A^0, A^1\} \subset \MN{2}$ defined by
    \[
        A^0 = \ketbra{0}{0} = \diag(1,0),
        \qquad
        A^1 = \ketbra{1}{1} = \diag(0,1).
    \]
    For a system of $N$ sites the resulting MPV is the GHZ state
    \[
        \ket{\mathrm{GHZ}_N}
        = \ket{0}^{\otimes N} + \ket{1}^{\otimes N}.
    \]
\end{definition}

\begin{theorem}[GHZ transfer map is idempotent]\label{thm:ghz_transfer_idempotent}
    \lean{MPSTensor.ghz_transferMap_idempotent}
    \leanok
    \uses{def:ghz_tensor, def:transfer_map}
    The transfer map of the GHZ tensor satisfies $\E_A^2 = \E_A$.
\end{theorem}

\begin{proof}
    \leanok
    The transfer map acts by
    $\E_A(X)_{ij} = \delta_{ij} X_{ij}$
    (extraction of the diagonal).
    This operation is manifestly idempotent.
\end{proof}

\begin{theorem}[GHZ tensor is an RFP]\label{thm:ghz_is_rfp}
    \lean{MPSTensor.ghz_isRFP}
    \leanok
    \uses{def:ghz_tensor, def:is_rfp, thm:ghz_transfer_idempotent}
    The GHZ tensor is a renormalization fixed point.
\end{theorem}

\begin{proof}
    \leanok
    \uses{thm:ghz_transfer_idempotent}
    Immediate from Theorem~\ref{thm:ghz_transfer_idempotent}.
\end{proof}

\begin{theorem}[GHZ tensor is not injective]\label{thm:ghz_not_injective}
    \lean{MPSTensor.ghz_not_isInjective}
    \leanok
    \uses{def:ghz_tensor, def:injective}
    The GHZ tensor is not injective: $\spn_{\C}\{A^0, A^1\}$ is the
    two-dimensional space of diagonal matrices in~$\MN{2}$, which does
    not equal~$\MN{2}$.
\end{theorem}

\begin{proof}
    \leanok
    Every element of $\spn_{\C}\{A^0, A^1\}$ has zero off-diagonal
    entries, so the matrix with all entries equal to~$1$ is not in the
    span.
\end{proof}

\begin{theorem}[$\Z_2$ on-site symmetry of the GHZ tensor]\label{thm:ghz_z2_symmetric}
    \lean{MPSTensor.ghz_isOnSiteSymmetric_Z2}
    \leanok
    \uses{def:ghz_tensor, def:on_site_symmetric}
    Let $\Z_2 = \{1, \sigma_x\}$ act on $\C^2$ by the Pauli-$X$
    matrix~$\sigma_x$.  The GHZ tensor is on-site symmetric under this
    action: for each $g \in \Z_2$, the twisted tensor is gauge
    equivalent to $A$.
\end{theorem}

\begin{proof}
    \leanok
    The identity element is trivial.  For the generator
    $g = \sigma_x$, the twisted tensor satisfies
    $(A^g)^i = A^{1-i}$, i.e.\ the two matrices are swapped.
    Conjugation by $\sigma_x$ implements the same swap, giving
    the required gauge equivalence.
\end{proof}

%% ===================================================================
\section{The AKLT state}\label{sec:aklt}
%% ===================================================================

The Affleck--Kennedy--Lieb--Tasaki (AKLT) state is the canonical
example of an injective MPS with non-trivial symmetry-protected
topological (SPT) order.

\begin{definition}[AKLT tensor]\label{def:aklt_tensor}
    \notready
    \uses{def:mps_tensor}
    Set $d = 3$ (spin-1) and $D = 2$.  The AKLT tensor is the family
    $\{A^0, A^1, A^2\} \subset \MN{2}$ defined via the Pauli matrices by
    \[
        A^0 = \frac{1}{\sqrt{2}}\,\sigma_z,
        \qquad
        A^1 = \frac{1}{\sqrt{2}}\,\sigma_+,
        \qquad
        A^2 = \frac{1}{\sqrt{2}}\,\sigma_-,
    \]
    where
    $\sigma_z = \ketbra{0}{0} - \ketbra{1}{1}$,
    $\sigma_+ = \sqrt{2}\,\ketbra{0}{1}$, and
    $\sigma_- = \sqrt{2}\,\ketbra{1}{0}$.
    Equivalently, the matrices are the Clebsch--Gordan coefficients
    coupling spin-$1$ and spin-$\tfrac12$ to spin-$\tfrac12$.
\end{definition}

\begin{theorem}[AKLT tensor is injective]\label{thm:aklt_injective}
    \notready
    \uses{def:aklt_tensor, def:injective}
    The AKLT tensor is injective:
    $\spn_{\C}\{A^0, A^1, A^2\} = \MN{2}$.
\end{theorem}

\begin{theorem}[AKLT tensor is normal]\label{thm:aklt_normal}
    \notready
    \uses{def:aklt_tensor}
    The transfer map of the AKLT tensor has a unique eigenvalue of
    modulus~$1$; in particular the AKLT tensor is normal.
\end{theorem}

\begin{theorem}[SU(2) on-site symmetry of the AKLT tensor]\label{thm:aklt_su2_symmetric}
    \notready
    \uses{def:aklt_tensor, def:on_site_symmetric}
    The AKLT tensor is on-site symmetric under the spin-$1$
    representation of $\mathrm{SU}(2)$, with the virtual
    representation being the fundamental (spin-$\tfrac12$)
    representation.  The projective phase is the non-trivial
    element of $H^2(\mathrm{SO}(3),\mathrm{U}(1)) \cong \Z_2$,
    witnessing a non-trivial SPT phase.
\end{theorem}

%% ===================================================================
\section{The cluster state}\label{sec:cluster}
%% ===================================================================

The cluster state is a universal resource for measurement-based quantum
computation and provides another example of a non-trivial SPT phase.

\begin{definition}[Cluster-state tensor]\label{def:cluster_tensor}
    \notready
    \uses{def:mps_tensor}
    Set $d = D = 2$.  The cluster-state tensor is the family
    $\{A^0, A^1\} \subset \MN{2}$ defined by
    \[
        A^0 = \frac{1}{\sqrt{2}}\,\Id_2,
        \qquad
        A^1 = \frac{1}{\sqrt{2}}\,\sigma_x,
    \]
    where $\sigma_x = \ketbra{0}{1} + \ketbra{1}{0}$.
\end{definition}

\begin{theorem}[Cluster-state tensor is not injective]\label{thm:cluster_not_injective}
    \notready
    \uses{def:cluster_tensor, def:injective}
    The cluster-state tensor is not injective:
    $\spn_{\C}\{A^0, A^1\} = \spn_{\C}\{\Id_2, \sigma_x\}$
    has dimension~$2$, not $\dim \MN{2} = 4$.
\end{theorem}

\begin{proof}
    \uses{def:cluster_tensor}
    Since $\sigma_x^2 = \Id_2$, the algebra generated by
    $\{\Id_2, \sigma_x\}$ is abelian and two-dimensional.
    Any matrix with distinct diagonal entries (e.g.\ $\sigma_z$)
    lies outside the span.
\end{proof}

\begin{remark}\label{rem:cluster_injective_blocked}
    Although the cluster-state tensor is not injective at length~$1$,
    it becomes injective after blocking with a suitable basis
    rotation.  Specifically, if one conjugates by the Hadamard
    matrix $H = \tfrac{1}{\sqrt{2}}(\sigma_x + \sigma_z)$
    to obtain $\tilde{A}^i = H A^i H^\dagger$, the resulting
    tensor has $\tilde{A}^0 = \tfrac{1}{\sqrt{2}} H$ and
    $\tilde{A}^1 = \tfrac{1}{\sqrt{2}} H \sigma_z$.
    The blocked tensor at length~$2$ in this basis is injective
    ($\dim \spn_{\C}\{\tilde{A}^w : |w| = 2\} = 4$).
\end{remark}

\begin{theorem}[$\Z_2$ on-site symmetry of the cluster-state tensor]\label{thm:cluster_z2_symmetric}
    \notready
    \uses{def:cluster_tensor, def:on_site_symmetric}
    Let $\Z_2 = \{1, \sigma_z\}$ act on $\C^2$ by the Pauli-$Z$
    matrix.  The cluster-state tensor is on-site symmetric under this
    action, with virtual representation given by $\sigma_z$.
    The projective phase is the non-trivial element of
    $H^2(\Z_2, \mathrm{U}(1)) \cong \Z_2$, witnessing a non-trivial
    SPT phase.
\end{theorem}
